\hypertarget{repository}{%
\section{Repository}\label{repository}}

All code written for this study is archived at https://github.com/hou1020/Parking. The directories referenced elsewhere in this dissertation are:

\par\addvspace{\medskipamount}\noindent
\begin{longtable}[]{@{}
  >{\raggedright\arraybackslash}p{(\columnwidth - 2\tabcolsep) * \real{0.3154}}
  >{\raggedright\arraybackslash}p{(\columnwidth - 2\tabcolsep) * \real{0.6646}}@{}}
\toprule
\begin{minipage}[b]{\linewidth}\raggedright
Directory
\end{minipage} &
\begin{minipage}[b]{\linewidth}\raggedright
Contents
\end{minipage} \\
\midrule
\endhead
\texttt{calculate/\allowbreak{}}
& Agreement against
the manual and OSM
references, polygon
filtering and result
merging \\
\texttt{analysis/\allowbreak{}} &
Validation, error
attribution,
sampling, ablation
and calibration
(§3.4--§3.9) \\
\texttt{fine-\allowbreak{}tuning/\allowbreak{}*.\allowbreak{}gpkg}
& The manual
reference labels and
the 1 km² grid \\
\texttt{fine-\allowbreak{}tuning/\allowbreak{}}
& Generic
fine-tuning of the
released checkpoint
(Appendix C) \\
\texttt{targeted-\allowbreak{}finetuning/\allowbreak{}}
& Targeted loss
weighting and the
threshold sweep
(Appendix C) \\
\texttt{parking-\allowbreak{}lot-\allowbreak{}mapping-\allowbreak{}tool/\allowbreak{}}
& The released
pipeline, with the
UK-specific tiling,
inference and
post-processing
written for this
study (§3.3) \\
\bottomrule
\end{longtable}
\par\addvspace{\medskipamount}\restoreparskip

\hypertarget{imagery}{%
\section{Imagery}\label{imagery}}

Getmapping aerial photography supplied through Digimap: 100 tiles at 0.25 m ground sample distance, three visible bands, EPSG:27700. The tile identifiers and version suffixes needed to reorder the same coverage are recorded in \texttt{parking-\allowbreak{}lot-\allowbreak{}mapping-\allowbreak{}tool/\allowbreak{}output\_\allowbreak{}files/\allowbreak{}tif\_\allowbreak{}processing\_\allowbreak{}progress.\allowbreak{}csv}, and every processing step from the raw tiles onward is reproducible from the code once the imagery is obtained under an equivalent licence.

\hypertarget{reference-data}{%
\section{Reference data}\label{reference-data}}

\begin{table}[htbp]
\centering
\begin{adjustbox}{max width=\textwidth}
\begin{tabular}{@{}lll@{}}
\toprule
Source & Use & Retrieved \\
\midrule
OpenStreetMap building footprints, road centrelines & Post-processing inputs (§3.3) & 25 June 2026 \\
OpenStreetMap land use, brownfield, pitch, \texttt{amenity=parking} & Error attribution only (§4.2) & 25 June 2026 \\
Ordnance Survey Open Greenspace & Sports facilities in error attribution & --- \\
\bottomrule
\end{tabular}
\end{adjustbox}
\end{table}
\restoreparskip

OpenStreetMap data are © OpenStreetMap contributors, available under the Open Database Licence. Ordnance Survey Open Greenspace is published under the Open Government Licence. Neither is used as ground truth; the distinction is set out in §3.1.

\hypertarget{reference-labels}{%
\section{Reference labels}\label{reference-labels}}

The 2,037 manually labelled car parks are held in the repository as GeoPackage, together with the 1 km² validation grid and the confidence attribute described in Appendix A. These are the labels against which every accuracy figure in Chapter 4 is measured, and they are original to this study.

\hypertarget{model}{%
\section{Model}\label{model}}

The segmentation network is the published checkpoint of Qiam, Devunuri and Lehe (2025), obtained from the authors' release and used without modification in the primary analysis. The fine-tuned checkpoints produced for Appendix C are derived works of that release and are not redistributed; the training code and logs that generate them are in \texttt{fine-\allowbreak{}tuning/\allowbreak{}} for the generic arm and \texttt{targeted-\allowbreak{}finetuning/\allowbreak{}} for the targeted loss weighting and threshold sweep.

\hypertarget{result-files}{%
\section{Result files}\label{result-files}}

Each appendix table is generated from a file in the repository rather than transcribed:

\par\addvspace{\medskipamount}\noindent
\begin{longtable}[]{@{}
  >{\raggedright\arraybackslash}p{(\columnwidth - 2\tabcolsep) * \real{0.2930}}
  >{\raggedright\arraybackslash}p{(\columnwidth - 2\tabcolsep) * \real{0.6870}}@{}}
\toprule
\begin{minipage}[b]{\linewidth}\raggedright
Table
\end{minipage} &
\begin{minipage}[b]{\linewidth}\raggedright
Source
\end{minipage} \\
\midrule
\endhead
Appendix B.1 &
\texttt{analysis/validation\_\allowbreak{}summary.csv} \\
Appendix B.2--B.4 &
\texttt{analysis/accuracy\_\allowbreak{}vs\_\allowbreak{}distance.csv},
\texttt{analysis/accuracy\_\allowbreak{}vs\_\allowbreak{}distance\_\allowbreak{}summary.csv} \\
Appendix C.1--C.2 &
\texttt{targeted-finetuning/Parking\_\allowbreak{}targeted\_\allowbreak{}run2/evaluation\_\allowbreak{}arms.csv} \\
Appendix C.3 &
\texttt{targeted-finetuning/Parking\_\allowbreak{}targeted\_\allowbreak{}run2/boundary\_\allowbreak{}bands\_\allowbreak{}arms.csv} \\
Appendix C.4 &
\texttt{targeted-finetuning/Parking\_\allowbreak{}targeted\_\allowbreak{}run2/selectivity.csv},
\texttt{standalone\_\allowbreak{}fp\_\allowbreak{}by\_\allowbreak{}category.csv} \\
Appendix C.5 &
\texttt{targeted-finetuning/Parking\_\allowbreak{}targeted\_\allowbreak{}run2/threshold\_\allowbreak{}sweep/generic\_\allowbreak{}threshold\_\allowbreak{}selected.csv} \\
\bottomrule
\end{longtable}
\par\addvspace{\medskipamount}\restoreparskip

\hypertarget{reproduction}{%
\section{Reproduction}\label{reproduction}}

The city-wide inference and the fine-tuning experiments require a GPU and were run in Google Colab; the notebooks (\texttt{run\_\allowbreak{}finetuning\_\allowbreak{}colab.\allowbreak{}ipynb}, \texttt{run\_\allowbreak{}targeted\_\allowbreak{}colab.\allowbreak{}ipynb}, \texttt{threshold\_\allowbreak{}sweep\_\allowbreak{}colab.\allowbreak{}ipynb}) pin \texttt{transformers==4.57.1}, otherwise running against the Colab environment's own package versions, and cache intermediate outputs, so a run interrupted partway resumes rather than restarting. All other analysis runs on CPU from the committed CSVs.
